\documentclass[9pt,a4paper]{article}
\usepackage[utf8]{inputenc}
\usepackage[T1]{fontenc}
\usepackage[french]{babel}
\usepackage{amsmath,amssymb,amsfonts}
\usepackage[portrait,left=1.5cm,right=1.5cm,top=1.2cm,bottom=0.8cm]{geometry}
\usepackage[dvipsnames]{xcolor}
\usepackage{multirow,tabularx,booktabs,array}
\usepackage{colortbl}
\usepackage{fouriernc}
\usepackage{multicol}
\usepackage{fancyhdr}

\DeclareMathOperator{\e}{e}

\pagestyle{fancy}
\renewcommand\headrulewidth{1pt}
\setlength{\headheight}{13.6pt}
\addtolength{\topmargin}{-1.6pt}
\fancyhead[L]{Lycée Denis Diderot}
\fancyhead[R]{Classe 1BTS}
\fancyfoot[C]{Grille de correction -- Mini Contrôle 3.3 -- Parité et Périodicité}
\fancyfoot[R]{Année 2025 -- 2026}
\fancyfoot[L]{Alexis RUIZ}

\definecolor{exohdr}{RGB}{30,100,180}
\definecolor{sousligne}{RGB}{220,232,248}
\definecolor{soustotal}{RGB}{255,220,80}

\renewcommand{\arraystretch}{1.3}

\newcommand{\exoheader}[3]{%
  \multicolumn{3}{|l|}{\cellcolor{exohdr}%
    \textcolor{white}{\textbf{Exercice #1 : #2}}\hfill
    \textcolor{white}{\textbf{/#3 pts}}%
  }\\
}
\newcommand{\colheader}{%
  \rowcolor{sousligne}
  \textbf{G} & \textbf{Réponse attendue} & \textbf{Pts} \\
}
\newcommand{\colheaderF}{%
  \rowcolor{sousligne}
  \textbf{F} & \textbf{Réponse attendue} & \textbf{Pts} \\
}
\newcommand{\soustotalrow}[2]{%
  \rowcolor{soustotal}
  \multicolumn{2}{|r|}{\textbf{Sous-total Exercice #1}} & \textbf{/#2} \\
}

\begin{document}

\begin{center}
  {\LARGE\textbf{GRILLE DE CORRECTION}}\\[3pt]
  Mini Contrôle 3.3 -- Parité et Périodicité \hspace{1cm} Classe 1BTS \hspace{1cm} \textbf{Total : /20 points}
\end{center}

\vspace{2mm}

%-----------------------------------------------------------------------
% EXERCICE 1
%-----------------------------------------------------------------------
\noindent
\begin{tabularx}{\linewidth}{|>{\centering\arraybackslash}p{0.55cm}|X|>{\centering\arraybackslash}p{1.1cm}|}
\hline
\exoheader{1}{Lecture graphique -- parité et périodicité}{8}
\hline
\colheader
\hline
1 & $\sin(\pi x)$ : \textbf{impaire} + \textbf{périodique $T=2$} & 1 \\
\hline
2 & $\cos(1{,}5\pi x)$ : \textbf{paire} + \textbf{périodique $T=\dfrac{4}{3}$} & 1 \\
\hline
3 & Signal triangulaire : \textbf{paire} + \textbf{périodique $T=4$} & 1 \\
\hline
4 & Signal en escalier : \textbf{périodique $T=3$} & 1 \\
\hline
5 & $\sin(\pi(x+1{,}3))$ : \textbf{périodique $T=2$} (ni paire, ni impaire) & 1 \\
\hline
6 & $0{,}5(\cos(1{,}5\pi x)-0{,}5)$ : \textbf{paire} + \textbf{périodique $T=\dfrac{4}{3}$} & 1 \\
\hline
7 & $0{,}32x\cos(\pi x)$ : \textbf{impaire} (non périodique) & 1 \\
\hline
8 & $0{,}1\sin(\pi x)\e^x$ : \textbf{Aucune réponse possible} & 1 \\
\hline
\soustotalrow{1}{8}
\hline
\end{tabularx}

{\footnotesize\textit{Barème ex.1 : 1 pt par graphique (tout ou rien). Une case incorrectement cochée annule le point.}}

\vspace{4mm}

%-----------------------------------------------------------------------
% EXERCICE 2
%-----------------------------------------------------------------------
\noindent
\begin{tabularx}{\linewidth}{|>{\centering\arraybackslash}p{0.55cm}|X|>{\centering\arraybackslash}p{1.1cm}|}
\hline
\exoheader{2}{Étude algébrique de la parité}{12}
\hline
\colheaderF
\hline
\multirow{2}{*}{$f$}
  & $f(-x)=3(-x)^2-1=3x^2-1=f(x)$ & 1 \\
\cline{2-3}
  & \textbf{Conclusion : $f$ est paire} & 1 \\
\hline
\multirow{2}{*}{$g$}
  & $g(-x)=(-x)^5+(-x)^3+(-x)=-x^5-x^3-x=-g(x)$ & 1 \\
\cline{2-3}
  & \textbf{Conclusion : $g$ est impaire} & 1 \\
\hline
\multirow{2}{*}{$h$}
  & $h(-x)=\dfrac{-5(-x)}{2(-x)^2+4}=\dfrac{5x}{2x^2+4}=-h(x)$ & 1 \\
\cline{2-3}
  & \textbf{Conclusion : $h$ est impaire} & 1 \\
\hline
\multirow{2}{*}{$j$}
  & $j(-x)=3(-x)^2+\dfrac{1}{(-x)^2}=3x^2+\dfrac{1}{x^2}=j(x)$ & 1 \\
\cline{2-3}
  & \textbf{Conclusion : $j$ est paire} & 1 \\
\hline
\multirow{2}{*}{$k$}
  & $k(-x)=(-x)^4-2(-x)^2+5=x^4-2x^2+5=k(x)$ & 1 \\
\cline{2-3}
  & \textbf{Conclusion : $k$ est paire} & 1 \\
\hline
\multirow{2}{*}{$l$}
  & $l(-x)=\dfrac{(-x)^3}{(-x)^2+1}=\dfrac{-x^3}{x^2+1}=-l(x)$ & 1 \\
\cline{2-3}
  & \textbf{Conclusion : $l$ est impaire} & 1 \\
\hline
\soustotalrow{2}{12}
\hline
\end{tabularx}

{\footnotesize\textit{Barème ex.2 : 1 pt calcul de $f(-x)$ + 1 pt conclusion. Vérifier que le domaine est mentionné comme symétrique.}}

\end{document}
